\subsection{Invariant factors of a submodule}

THEOREM 1. --- Let L be a free module over a principal ideal domain A, and let M be a submodule of L of finite rank n.~Then there exist a basis B of L, n elements \(e_{i}\) of B, and n nonzero elements \(\alpha_{i}\) of A ( \(1 \leqslant i \leqslant n\) ) such that :

\begin{enumerate}
\def\labelenumi{\alph{enumi})}
\item
  the \(\alpha_{i}e_{i}\) form a basis of \(\mathbf{M}\) ;
\item
  \(\alpha_{i}\) divides \(\alpha_{i + 1}\) for \(1\leqslant i\leqslant n - 1\)
\end{enumerate}

Moreover the module M' generated by the \((e_{i})\) and the principal ideals \(A\alpha_{i}\) are uniquely determined by the above conditions; the module M'/M is the torsion submodule of L/M, and is isomorphic to the direct sum of the A-modules A/A \(\alpha_{i}\) ; finally L/M is the direct sum of M'/M and a free module isomorphic to L/M'.

\begin{enumerate}
\def\labelenumi{\arabic{enumi})}
\tightlist
\item
  Existence of the \(e_i\) and the \(\alpha_i\) .
\end{enumerate}

If \(M = \{0\}\) then the theorem is trivial. If \(M \neq \{0\}\) then it follows from Prop. 3 that there exist an element \(e_{1}\) of L, a nonzero element \(\alpha_{1}\) of A, and a submodule \(L_{1}\) of L, such that L is the direct sum of \(Ae_{1}\) and \(L_{1}\) , such that M is the direct sum of \(A\alpha_{1}e_{1}\) and the submodule \(M_{1} = M \cap L_{1}\) of L, and such that \(g(M) \subset A\alpha_{1}\) for every linear form g on L.

Now we can proceed by induction on the rank \(n\) of \(M\) . Since \(L_1\) is a free module (VII, p.~15, Cor. 2) and \(M_1\) has rank \(n - 1\) , there exist a basis \(B_1\) of \(L_1\) , \((n - 1)\) elements \(e_2, \ldots, e_n\) of \(B_1\) , and nonzero elements \(\alpha_2, \ldots, \alpha_n\) of \(A\) such that \((\alpha_2e_2, \ldots, \alpha_ne_n)\) is a basis of \(M_1\) , and \(\alpha_i\) divides \(\alpha_{i+1}\) for \(2 \leqslant i \leqslant n - 1\) . If \(L'\) is the submodule of \(L_1\) generated by the elements of \(B_1\) distinct from \(e_2, \ldots, e_n\) , then \(L\) is the direct sum of \(L'\) and the module \(M'\) generated by \(e_1, \ldots, e_n\) ; now \((e_1, \ldots, e_n)\) is a basis for \(M'\) , and \((\alpha_1e_1, \ldots, \alpha_ne_n)\) is a basis for \(M\) . It now remains only to show that \(\alpha_1\) divides \(\alpha_2\) ; but \(A\alpha_2\) has the form \(g(M_1)\) , where \(g\) is the linear form on \(L\) defined by \(g(e_2) = 1\) , \(g(e_i) = 0\) for \(i \neq 2\) , and \(g(L') = \{0\}\) ; and we have seen above that \(g(M_1) \subset A\alpha_1\) .

2) Uniqueness properties.

Since the \(\alpha_{i}\) are nonzero, the module \(\mathbf{M}'\) is the set of \(x\in L\) such that \(\beta x\in M\) for some \(\beta \neq 0\) in A; in other words \(\mathbf{M}' / \mathbf{M}\) is the torsion submodule of L/M. This uniquely determines \(\mathbf{M}'\) .

It is clear that \(M^{\prime}/M\) is isomorphic to the direct sum of the n cyclic modules \(A/A\alpha_{i}\) (II, p.~204, formula (26)). Let r be the number of the ideals \(A\alpha_{i}\) which are distinct from A: so that the first n-r ideals \(A\alpha_{i}\) are equal to A, and the last r are distinct from A. Then \(M^{\prime}/M\) is also isomorphic to the direct sum of the modules \(A/A\alpha_{n}, \ldots, A/A\alpha_{n-r+1}\) , where \(A\alpha_{n} \subset A\alpha_{n-1} \subset \ldots \subset A\alpha_{n-r+1} \neq A\) .

Thus the conditions of the Cor. to Prop. 2 (VII, p.~16) are satisfied: the ideals \(\mathbf{A}\alpha_{i}\) ( \(1 \leqslant i \leqslant n\) ) are thus uniquely determined by the module \(\mathbf{M}'/\mathbf{M}\) .

Since L is the direct sum of \(M'\) and \(L'\) , it follows that L/M is the sum of \(M'/M\) and \((L'+M)/M\) , and this sum is direct since \((L'+M) \cap M' = M\) ; on the other hand \((L'+M)/M\) is isomorphic to \(L'/(M \cap L')\) (I, p.~41, Th. 4, c)), that is to \(L'\) , which shows that \((L'+M)/M\) is a free module isomorphic to \(L/M'\) .

COROLLARY. --- A submodule M of finite rank in a free module L over a principal ideal domain A has a complement in L if and only if L/M is torsion free.

In the notation of Th. 1, if L/M is torsion free, then \(M = M'\) , and \(M'\) has a complement \(L'\) in L. Conversely if M has a complement \(L'\) in L, then L/M is isomorphic to \(L'\) , which is free (VII, p.~15, Cor. 2), and a fortiori torsion free.

Remark. --- It can happen that a submodule M of infinite rank in a free module L can be such that L/M is torsion free, but M has no complement in L (VII, p.~60, Ex. 6, b)).

DEFINITION 1. --- In the notation and hypotheses of Th. 1, the ideals \(A\alpha_{i}\) of A are called the invariant factors of the submodule M with respect to the module L.

In the case where A is either the ring Z of rational integers or the polynomial ring K{[}X{]} in one indeterminate over a field K, there is a canonical way of choosing a generator for each ideal of A: a positive integer in the case of Z, or a monic polynomial in the case of K{[}X{]} (VII, p.~5). In each of these cases, the canonical generator of the invariant factor \(A\alpha_{i}\) is also called an invariant factor of M with respect to L, by abuse of language.

